% Типовой расчёт № 2 по линейной алгебре
% Лаборатория математики ФН1
% Компиляция: pdflatex tr-02-usloviya.tex

\documentclass[12pt, a4paper]{article}

\usepackage[T2A]{fontenc}
\usepackage[utf8]{inputenc}
\usepackage[russian]{babel}
\usepackage[left=25mm, right=20mm, top=20mm, bottom=20mm]{geometry}
\usepackage{amsmath, amssymb}
\usepackage{enumitem}

\pagestyle{plain}
\setlength{\parindent}{1.25cm}

\begin{document}

\begin{center}
  {\large\bfseries Типовой расчёт № 2}

  \vspace{0.3em}
  {\large Линейные пространства и операторы}

  \vspace{0.5em}
  {\small Линейная алгебра и аналитическая геометрия · 1 курс, 1 семестр\\
  Лаборатория математики ФН1}
\end{center}

\vspace{1em}

\noindent
Числовые данные вариантов — в файле \texttt{tr-02-varianty.csv}.

\section*{Задача 1. Линейная зависимость}

Выяснить, образует ли система векторов базис пространства $\mathbb{R}^3$:
\[
e_1 = (1,\ a,\ 2),\qquad e_2 = (0,\ 1,\ -1),\qquad e_3 = (a,\ 2,\ 1).
\]
При каких значениях $a$ система линейно зависима? Для одного из таких значений
указать нетривиальную линейную комбинацию, равную нулю.

\section*{Задача 2. Координаты в базисе}

Разложить вектор $x = (b,\ 1,\ 3)$ по базису из задачи~1 при том значении $a$,
для которого базис существует. Записать матрицу перехода от стандартного
базиса к данному.

\section*{Задача 3. Матрица линейного оператора}

Линейный оператор $\varphi$ действует по правилу
\[
\varphi(x_1,\ x_2,\ x_3) = (x_1 + c\,x_2,\ \ x_2 - x_3,\ \ c\,x_1 + x_3).
\]
Записать матрицу оператора в стандартном базисе, найти ядро и образ,
проверить равенство $\dim\ker\varphi + \operatorname{rk}\varphi = 3$.

\section*{Задача 4. Собственные значения и векторы}

Найти собственные значения и собственные векторы матрицы
\[
A = \begin{pmatrix} d & 1 & 0 \\ 1 & d & 1 \\ 0 & 1 & d \end{pmatrix}.
\]
Выяснить, диагонализуема ли матрица; если да~--- записать разложение
$A = S\Lambda S^{-1}$.

\section*{Задача 5. Квадратичная форма}

Привести квадратичную форму к каноническому виду методом Лагранжа
и определить её знакоопределённость:
\[
Q(x_1, x_2, x_3) = x_1^2 + 2k\,x_1x_2 + 2x_2^2 - x_3^2 + 2x_2x_3.
\]

\section*{Задача 6. Ортогонализация}

Применить процесс Грама~--- Шмидта к системе векторов из задачи~1
и получить ортонормированный базис. Проверить ортогональность
попарными скалярными произведениями.

\vspace{1.2em}

\section*{Требования к оформлению}

Те же, что и для типового расчёта №~1. Для задач~4 и~6 обязательна проверка:
подстановка собственных векторов в равенство $Ax = \lambda x$ и вычисление
скалярных произведений соответственно.

\section*{Критерии оценивания}

\begin{center}
\begin{tabular}{|c|l|}
  \hline
  Баллы & За что начисляются \\
  \hline
  2 & задача решена верно, преобразования обоснованы \\
  1 & верный метод, допущена вычислительная ошибка \\
  0 & метод выбран неверно либо решение отсутствует \\
  \hline
\end{tabular}
\end{center}

\vspace{0.5em}
\noindent
Расчёт зачтён при \textbf{8 баллах из 12} и успешной защите.

\vspace{0.7em}
\noindent\textbf{Сроки.} Выдача~--- 10-я учебная неделя,
сдача~--- \textbf{до конца 15-й недели}.

\end{document}
